\section*{References}
\begingroup
\setlength{\parindent}{0pt}\setlength{\parskip}{5pt}
\newcommand{\entry}[1]{\par\hangindent=1.5em\hangafter=1 #1}

\entry{Barnes, L.~A. (2012). The fine-tuning of the universe for intelligent life. \emph{Publications of the Astronomical Society of Australia}, 29(4), 529--564.}
\entry{Bostrom, N. (2003). Are you living in a computer simulation? \emph{Philosophical Quarterly}, 53(211), 243--255.}
\entry{Bourget, D., and Chalmers, D.~J. (2014). What do philosophers believe? \emph{Philosophical Studies}, 170(3), 465--500.}
% VERIFY: exact title/volume
\entry{Bystranowski, P., Dranseika, V., and \.{Z}uradzki, T. (2022). Half a century of bioethics and philosophy of medicine: A topic-modeling study. \emph{Bioethics}, 36(9), 902--925.}
\entry{Cahn, S.~M. (1977). Cacodaemony. \emph{Analysis}, 37(1), 69--73.}
% VERIFY: which Adams work is cited (in-text: Adams 2019)
\entry{Adams, [initial]. (2019). [Title to be verified against the in-text citation.]}
% VERIFY: Allison publication details
\entry{Allison, D.~C. (2021). \emph{The Resurrection of Jesus: Apologetics, Polemics, History}. T\&T Clark.}
% VERIFY: De Cruz 2017 exact study
\entry{De Cruz, H. (2017). Religious disagreement: An empirical study among academic philosophers. \emph{Episteme}, 14(1), 71--87.}
% VERIFY: De Cruz 2018 exact study
\entry{De Cruz, H. (2018). [Qualitative study of religious belief trajectories among philosophers of religion; verify exact work.]}
% VERIFY: year (book is commonly dated 2015)
\entry{De Cruz, H., and De Smedt, J. (2016). \emph{A Natural History of Natural Theology: The Cognitive Science of Theology and Philosophy of Religion}. MIT Press.}
\entry{Draper, P. (1989). Pain and pleasure: An evidential problem for theists. \emph{No\^{u}s}, 23(3), 331--350.}
\entry{Draper, P., and Nichols, R. (2013). Diagnosing bias in the philosophy of religion. \emph{The Monist}, 96(3), 420--446.}
% VERIFY: Gauch 2013 venue
\entry{Gauch, H.~G., Jr. (2013). The methodology of ramified natural theology. \emph{Philosophia Christi}, 15(2), 283--298.}
\entry{Geach, P.~T. (1973). Omnipotence. \emph{Philosophy}, 48(183), 7--20.}
\entry{Hartshorne, C. (1984). \emph{Omnipotence and Other Theological Mistakes}. State University of New York Press.}
\entry{Hasker, W. (1989). \emph{God, Time, and Knowledge}. Cornell University Press.}
% VERIFY: Holder 2021 details
\entry{Holder, R. (2021). \emph{Ramified Natural Theology in Science and Religion: Moving Forward from Natural Theology}. Routledge.}
% VERIFY: Lapide edition/year
\entry{Lapide, P. (1983). \emph{The Resurrection of Jesus: A Jewish Perspective}. Augsburg.}
\entry{Lessing, G.~E. (1777). \"{U}ber den Beweis des Geistes und der Kraft. English translation: On the proof of the spirit and of power. In H.~Chadwick (ed.), \emph{Lessing's Theological Writings}. Stanford University Press, 1957.}
\entry{L\"{u}demann, G. (1994). \emph{The Resurrection of Jesus: History, Experience, Theology}. Fortress Press.}
\entry{McGrew, T. (2004). Has Plantinga refuted the historical argument? \emph{Philosophia Christi}, 6(1), 7--26.}
% VERIFY: McGrew and McGrew 2006 details
\entry{McGrew, L., and McGrew, T. (2006). On the historical argument: A rejoinder to Plantinga. \emph{Philosophia Christi}, 8(1), 23--38.}
\entry{Mill, J.~S. (1874). Theism. In \emph{Three Essays on Religion}. Longmans, Green.}
\entry{Morris, T.~V. (1986). \emph{The Logic of God Incarnate}. Cornell University Press.}
\entry{Morriston, W. (2000). Must the beginning of the universe have a personal cause? \emph{Faith and Philosophy}, 17(2), 149--169.}
% VERIFY: which Morriston 2020 work
\entry{Morriston, W. (2020). [Title to be verified against the in-text citation.]}
% VERIFY: Nickel 2015 details
\entry{Nickel, [initial]. (2015). [Dwindling-probabilities exchange contribution; verify author and venue.]}
% VERIFY: Ocampo 2024 details
\entry{Ocampo, [initial]. (2024). [Systematic map of Stage~II bridging strategies; verify author and venue.]}
\entry{Oppy, G. (2006). \emph{Arguing about Gods}. Cambridge University Press.}
\entry{Plantinga, A. (1974). \emph{The Nature of Necessity}. Oxford University Press.}
\entry{Plantinga, A. (2000). \emph{Warranted Christian Belief}. Oxford University Press.}
\entry{Pruss, A.~R. (2009). The Leibnizian cosmological argument. In W.~L. Craig and J.~P. Moreland (eds.), \emph{The Blackwell Companion to Natural Theology} (pp. 24--100). Wiley-Blackwell.}
\entry{Pruss, A.~R., and Rasmussen, J.~L. (2018). \emph{Necessary Existence}. Oxford University Press.}
\entry{Rasmussen, J. (2009). From a necessary being to God. \emph{International Journal for Philosophy of Religion}, 66(1), 1--13.}
% VERIFY: in-text "Rea 2005" may be Brower and Rea 2005
\entry{Brower, J.~E., and Rea, M.~C. (2005). Material constitution and the Trinity. \emph{Faith and Philosophy}, 22(1), 57--76.}
\entry{Roche, W., and Shogenji, T. (2014). Dwindling confirmation. \emph{Philosophy of Science}, 81(1), 114--137.}
\entry{Rowe, W.~L. (1979). The problem of evil and some varieties of atheism. \emph{American Philosophical Quarterly}, 16(4), 335--341.}
% VERIFY: Russell brute-fact citation form (1948 BBC debate)
\entry{Russell, B., and Copleston, F.~C. (1948). The existence of God: A debate. Reprinted in B.~Russell, \emph{Why I Am Not a Christian}. Allen \& Unwin, 1957.}
\entry{Swinburne, R. (2003). \emph{The Resurrection of God Incarnate}. Oxford University Press.}
\entry{Swinburne, R. (2004). Natural theology, its ``dwindling probabilities'' and ``lack of rapport''. \emph{Faith and Philosophy}, 21(4), 533--546.}
\entry{Tegmark, M. (2008). The mathematical universe. \emph{Foundations of Physics}, 38(2), 101--150.}
\entry{Tuggy, D. (2003). The unfinished business of Trinitarian theorizing. \emph{Religious Studies}, 39(2), 165--183.}
% VERIFY: in-text "Vestrup 2001" may be McGrew, McGrew and Vestrup 2001
\entry{McGrew, T., McGrew, L., and Vestrup, E. (2001). Probabilities and the fine-tuning argument: A sceptical view. \emph{Mind}, 110(440), 1027--1037.}
\entry{White, R. (2000). Fine-tuning and multiple universes. \emph{No\^{u}s}, 34(2), 260--276.}
% VERIFY: Wildman 2004 details
\entry{Wildman, W.~J. (2004). The Divine Action Project, 1988--2003. \emph{Theology and Science}, 2(1), 31--75.}
% VERIFY: Wiles work cited
\entry{Wiles, M. (1986). \emph{God's Action in the World}. SCM Press.}
\entry{Wright, N.~T. (2003). \emph{The Resurrection of the Son of God}. Fortress Press.}
\entry{Wykstra, S.~J. (1984). The Humean obstacle to evidential arguments from suffering: On avoiding the evils of ``appearance''. \emph{International Journal for Philosophy of Religion}, 16(2), 73--93.}
\endgroup
